%! Interpreting p-Values — Unit 6, Lesson 6.5 — Guided Notes
\documentclass[10pt]{article}
\usepackage{apstats-article}
\usepackage{apstats-boxes}
\newcommand{\TallMath}[1]{$\displaystyle #1\rule[-1.4em]{0pt}{3.2em}$}

\begin{document}

\pageheader{Unit 6, Lesson 6.5}{Guided Notes}
\namedateperiod

\vspace{0.15in}

\begin{objectivebox}
Compute and interpret a $p$-value from a one-sample $z$-test statistic for a proportion.
\end{objectivebox}

\begin{vocabbox}
\begin{description}
  \item[\termblanklong{$p$-value}] The probability of observing a sample result as extreme as or
    more extreme than the one obtained, \emph{assuming \blank{2cm} is true}. A probability between
    \blank{1.5cm} and \blank{1.5cm}.
  \item[\termblanklong{``As extreme''}] In the direction(s) specified by \blank{3cm}.
\end{description}
\end{vocabbox}

\begin{notesbox}{Finding the $p$-Value: Which Tail?}
\renewcommand{\arraystretch}{2.0}
\begin{tabularx}{\linewidth}{@{} p{0.22\linewidth} X X @{}}
  $H_a$         & \textbf{``Extreme'' means} & \textbf{$p$-value formula} \\ \hline
  $p > p_0$     & $z$ is \blank{3cm}         & $P(Z \ge z) = 1 - \Phi(z)$ \\
  $p < p_0$     & $z$ is \blank{3cm}         & $P(Z \le z) = \Phi(z)$ \\
  $p \ne p_0$   & $|z|$ is \blank{3cm}       & $2P(Z \ge |z|)$ \\
\end{tabularx}
\end{notesbox}

\begin{notesbox}{Interpreting the $p$-Value}
\textbf{Template:} ``Assuming [restate $H_0$ in context], there is a \blank{2cm} probability
of observing \blank{6cm} (or more extreme) by chance alone.''

\vspace{4pt}
\textbf{Critical errors to avoid:}
\begin{itemize}
  \item[\textbf{$\times$}] ``The $p$-value is the probability that $H_0$ is \blank{2cm}''
  \item[\textbf{$\times$}] ``The $p$-value is the probability the result was due to \blank{2cm}''
  \item[\checkmark] ``Assuming $H_0$ is true, \blank{6cm}''
\end{itemize}
\end{notesbox}

\begin{notesbox}{Worked Example: Defective Widget Rate}
A manufacturer claims 10\% defect rate ($p_0=0.10$). A random sample of 150 widgets finds
$\hat p = 22/150 \approx 0.147$. Test $H_a: p \ne 0.10$ (two-sided). $z \approx 1.92$.

\begin{itemize}
  \item Direction: two-sided, so $p$-value $= 2P(Z \ge \blank{1.5cm}) = 2(1-\Phi(\blank{1.5cm})) = 2(\blank{2cm}) \approx$ \blank{2cm}
  \item Interpret: ``Assuming 10\% of widgets are defective, there is a \blank{2cm} probability
        of obtaining a sample proportion as far from 0.10 as 0.147 by \blank{3cm} alone.''
\end{itemize}
\end{notesbox}

\begin{notesbox}{What the $p$-Value Does NOT Tell You}
\begin{enumerate}
  \item The $p$-value is \textbf{not} the probability \blank{8cm} is true.
  \item A large $p$-value does \textbf{not} \blank{8cm} $H_0$.
  \item A small $p$-value does \textbf{not} mean the effect is \blank{4cm} significant.
\end{enumerate}
\end{notesbox}

\begin{practicebox}
For each scenario, identify the correct tail and compute the $p$-value.

\begin{enumerate}[label=\alph*.]
  \item $H_a: p > 0.60$, $z = 1.45$. Tail: \blank{2cm} \quad $p$-value $=$ \blank{3cm}
  \item $H_a: p < 0.30$, $z = -2.10$. Tail: \blank{2cm} \quad $p$-value $=$ \blank{3cm}
  \item $H_a: p \ne 0.50$, $z = -1.75$. Tail: \blank{2cm} \quad $p$-value $=$ \blank{3cm}
\end{enumerate}
\end{practicebox}

\end{document}
